\documentclass[11pt,a4paper]{article}

\usepackage[spanish,es-noshorthands]{babel}
\usepackage[T1]{fontenc}
\usepackage[utf8]{inputenc}
\usepackage[a4paper,margin=2.4cm]{geometry}
\usepackage{amsmath,amssymb}
\usepackage{listings}
\usepackage[colorlinks=true,linkcolor=blue,urlcolor=blue]{hyperref}

\lstset{
  basicstyle=\ttfamily\small,
  frame=single,
  breaklines=true,
  columns=fullflexible,
  keepspaces=true
}

\title{\textbf{Ejercicio Práctico 2}\\[2mm]
\large Modelado de un Registro con Invariante}
\author{Proyecto: Sistema de Planificación de Viajes en Transporte Público}
\date{\today}

\begin{document}

\maketitle

\section{Objetivo}

Modelar un \textbf{registro} (\emph{record}) para un componente del proyecto e indicar al menos una \textbf{invariante}.

\section{Definición del registro}

Un registro es un tipo compuesto por campos con nombre. Se modelan dos registros del subsistema de rutas: \texttt{Punto}, que representa una coordenada geográfica, y \texttt{Recorrido}, que agrupa los datos de un trayecto.

\begin{lstlisting}
Punto :: lat : real
         lng : real;

Recorrido :: ruta         : RutaId
             sentido      : Sentido
             coordenadas  : seq of Punto
             color        : seq of char;
\end{lstlisting}

\section{Invariante}

Se definen dos invariantes. La primera garantiza que las coordenadas estén dentro de un rango geográfico válido; la segunda, que todo recorrido tenga al menos dos puntos, coincidiendo con la regla de negocio del sistema (un trayecto requiere un punto de origen y uno de destino).

\begin{lstlisting}
inv punto ==
  lat >= -90  and lat <= 90  and
  lng >= -180 and lng <= 180;

inv recorrido ==
  len coordenadas >= 2;
\end{lstlisting}

\section{Justificación}

\begin{itemize}
  \item La invariante de \texttt{Punto} impide que el sistema almacene coordenadas geográficamente inválidas.
  \item La invariante de \texttt{Recorrido} garantiza que no existan recorridos vacíos o de un solo punto, una propiedad necesaria para que el algoritmo de búsqueda de rutas produzca resultados con sentido.
\end{itemize}

\end{document}
